\subsection{Distance globale entre séquences}
Pour calculer la distance de Levenshtein~\cite{levenshtein}, Wagner et Fischer~\cite{wagner} considèrent la dernière opération d’un alignement. Si \(D_{ij}\) est la distance entre les préfixes de longueurs \(i\) et \(j\), on obtient, pour des coûts unitaires,
\[
 D_{ij}=\min\{D_{i-1,j}+1,\ D_{i,j-1}+1,\ D_{i-1,j-1}+[A_i\ne B_j]\},
 \qquad D_{i0}=i,\quad D_{0j}=j.
\]
Les deux premiers termes correspondent à une délétion et à une insertion. Le troisième aligne deux lettres, au coût de zéro si elles coïncident et de un sinon. Une suite d'éditions qui mène de \(A\) à \(B\) peut être inversée, d'où la symétrie. L'écart de longueur donne la borne inférieure \(\lvert m-n\rvert\) et les substitutions suivies de l'ajustement de longueur donnent la borne supérieure \(\max(m,n)\).

Pour établir la récurrence de NvPD~\cite{sadiq}, notons \(P_j=D_{i-1,j}\) et \(k\) la dernière position de \(A_i\) dans \(B_{1:j}\), avec \(k=0\) en l’absence de correspondance. La démonstration repose sur l’inégalité \(P_v\le P_u+v-u\) pour \(v\ge u\), car des insertions suffisent pour allonger le préfixe de \(B\). Considérons la dernière opération qui consomme \(A_i\), suivie d'insertions. Si elle supprime \(A_i\) en colonne \(t\), son coût total est au moins \(P_j+1\). Si elle substitue une lettre en \(t\), il est au moins \(P_{j-1}+1\). Pour une correspondance, \(P_{t-1}+j-t\) décroît au sens large avec \(t\). Ce coût est donc minimal à la dernière occurrence, comme l’illustre la figure~\ref{fig:dep}. Si \(k=0\), l’algorithme utilise \(P_0+j=i-1+j\). Pour \(i,j\ge1\), cette valeur est au moins égale à \(\max(i,j)\) et laisse le minimum inchangé.
\schema{dependances}{Les flèches indiquent les prédécesseurs de la case terminale \(D_{1,2}\). La dernière occurrence de \texttt{A} dans \texttt{AC} se trouve en colonne \(k=1\). NvPD lit donc le score \(P_0\) de la ligne précédente pour atteindre \(j=2\).}{fig:dep}{src/cpu/dp.cpp}

Dans l’algorithme~\ref{alg:nvpd}, \(K\) conserve la dernière occurrence de chaque lettre à chaque colonne. Les tableaux \(P\) et \(Q\) échangent leurs rôles dès que tous les scores d’une ligne sont disponibles.
\begin{algorithm}[!htbp]
\caption{Mise à jour de NvPD par colonnes indépendantes}\label{alg:nvpd}
\begin{algorithmic}[1]
\Require Strings \(A[1:m], B[1:n]\)
\Ensure Unit-cost global edit distance \(d(A,B)\)
\State Build \(K[c,j]=\max(\{t\le j:B_t=c\}\cup\{0\})\) for symbols \(c\) in \(A\)
\State \(P[0:n]\gets (0,1,\ldots,n)\)
\For{\(i=1,\ldots,m\)}
  \State \(Q[0]\gets i\)
  \ForAll{\(j=1,\ldots,n\) in parallel}
    \State \(k\gets K[A_i,j]\)
    \State \(Q[j]\gets\min\{1+\min(P[j-1],P[j]),\,P[\max(0,k-1)]+j-k\}\)
  \EndFor
  \State Synchronize workers and swap \(P,Q\)
\EndFor
\State \Return \(P[n]\)
\end{algorithmic}
\end{algorithm}

\synthese{Le calcul du score ne conserve que deux lignes. Une barrière par ligne suffit pour que tous les fils accèdent aux valeurs dont dépend la suivante.}

\subsection{Architecture des noyaux d’alignement}
La disposition des prédécesseurs détermine la mémoire à conserver et la fréquence des synchronisations. La programmation dynamique (DP, \emph{dynamic programming}) séquentielle avance par lignes dans deux tableaux. Sur l'anti-diagonale \(i+j=s\), les valeurs d'entrée appartiennent aux niveaux \(s-1\) et \(s-2\). Trois tableaux tournants suffisent et chacune des \(m+n-1\) anti-diagonales intérieures comporte au plus \(\min(m,n)\) cases. NvPD avance en \(m\) lignes de \(n\) colonnes. Le mode coût seul place la chaîne courte sur les lignes. Toutes ces méthodes traitent \(mn\) cases.
\schema{architecture}{Les exécutables local et distribué lisent le FASTA. Le premier peut restituer un chemin optimal pour une paire. Le second distribue le calcul des scores entre les rangs MPI.}{fig:arch}{src/cli/main.cpp}
Le lecteur de la figure~\ref{fig:arch} fournit les mêmes chaînes à tous les noyaux. La bibliothèque C++ Kokkos, d’Edwards, Trott et Sunderland, exécute une même récurrence sur processeur ou sur GPU~\cite{kokkos}. Edlib~1.2.7, de Martin Šošić et Mile Šikić, adapte l’algorithme de Myers fondé sur les vecteurs de bits~\cite{edlib,myers,edlibcode}. WFA2-lib~2.3.6, de Santiago Marco-Sola et ses coauteurs, explore les alignements par fronts de coût croissant~\cite{wfa,wfa2code}. BiWFA propage un front depuis chaque extrémité afin de reconstruire le chemin en mémoire linéaire dans le score~\cite{biwfa}.

QuickEd~1.2, de Doblas et ses coauteurs, utilise des fenêtres pour estimer une borne puis calcule un alignement en bande~\cite{quicked,quickedcode}. A*PA2, de Groot Koerkamp, combine calcul en bits et borne inférieure sur le coût restant~\cite{astar,astarcode}. Les adaptateurs d’Edlib et de WFA2 fournissent le score global ou le chemin optimal. QuickEd dispose d'un exécutable séparé et reconstruit aussi l'alignement. Les pilotes A*PA2 et eWFA-GPU appellent les implémentations des auteurs. Accelign, de Kallenborn et ses coauteurs, calcule les scores par tuiles sur GPU et échange les prédécesseurs entre registres d'un warp~\cite{accelign}. Sa comparaison sur dix millions de lectures de 148 bases porte sur le score.

La comparaison des durées porte sur des alignements globaux. D’autres outils répondent à des questions voisines. Sassy, de Beeloo et Groot Koerkamp, recherche un motif approximatif dans un texte par calcul SIMD~\cite{sassy}. La recherche autorise un début et une fin libres dans le texte. Le motif \texttt{ACG} apparaît dans \texttt{TTACGTT} avec un coût de recherche nul. La distance globale de ces chaînes vaut quatre. À l’échelle d’un génome, mm2-plus parallélise le chaînage des ancres de minimap2 pour repérer les segments à aligner~\cite{mm2plus}.

Le mode trace des noyaux denses conserve \((m+1)(n+1)\) entiers. Il remonte depuis \((m,n)\) et écrit les opérations dans l'ordre de transformation de \(A\) en \(B\).\footnote{Un CIGAR indique la longueur et le type des opérations. Les tests le rejouent pour vérifier les bases consommées et le coût.}

\schema{execution}{Une équipe NvPD synchronise ses plages de colonnes à chaque ligne. Pour une collection, deux fils peuvent enchaîner des paires indépendantes selon les affectations illustrées.}{fig:execution}{src/cpu/dp.cpp}
\FloatBarrier
La figure~\ref{fig:execution} montre comment les mêmes cœurs peuvent se partager les colonnes d’une paire ou calculer simultanément plusieurs paires. Lorsque la commande \texttt{matrix} dispose d’au moins une paire par fil, elle affecte un seul travailleur à chaque appel de \texttt{nvpd-omp}.

\synthese{Pour seize séquences, les \(16\times15/2=120\) alignements du triangle supérieur peuvent s’exécuter indépendamment. À l’intérieur de chaque alignement, le travail dépend du domaine que l’algorithme doit explorer.}

\subsection{Recherche dans une bande de coût borné}
Une borne sur le nombre d'éditions permet de limiter le nombre de diagonales explorées. Chaque insertion ou délétion déplace le chemin d'une diagonale. Un chemin de coût au plus \(k\) reste donc dans \(\lvert i-j\rvert\le k\). Le programme double le seuil jusqu'au premier score admissible, suivant le principe du calcul en bande~\cite{ukkonen}. Pour \(L=\max(m,n)\) et une distance finale \(d\), la somme des largeurs donne \(O(L(d+1))\) opérations et \(O(d+1)\) cases de mémoire.
L’algorithme~\ref{alg:band} s’arrête dès que le score terminal est inférieur ou égal au seuil \(k\). Tout chemin de coût moindre tiendrait aussi dans cette bande.
\begin{algorithm}[!htbp]
\caption{Alignement global par élargissement de bande}\label{alg:band}
\begin{algorithmic}[1]
\Require Strings \(A,B\), optional common-end trimming
\Ensure Unit-cost global edit distance \(d(A,B)\)
\State Trim common ends if requested and return the other length if one string is empty
\State \(m\gets |A|,\; n\gets |B|,\; L\gets\max(m,n),\; k\gets\max(1,|m-n|)\)
\Loop
  \State Fill two DP rows over \(\max(0,i-k)\le j\le\min(n,i+k)\) using out-of-band sentinels
  \If{the terminal cell is reachable and \(D_{mn}\le k\)} \State \Return \(D_{mn}\) \EndIf
  \State \(k\gets\min(2k,L)\)
\EndLoop
\end{algorithmic}
\end{algorithm}
Avec des coûts unitaires, supprimer un préfixe ou un suffixe identique conserve \(d\). Dans le calcul bit-parallèle de Myers, un mot machine contient les différences de plusieurs scores successifs. WFA2 visite les diagonales qu’un front de faible coût peut atteindre. Ding et coll. accélèrent l’exploration par fronts avec des requêtes de plus long préfixe commun entre suffixes~\cite{ding}.

\synthese{À un budget de \(k\) éditions correspondent au plus \(2k+1\) diagonales par ligne. Un score terminal inférieur ou égal à ce budget suffit pour arrêter le calcul. L’exemple suivant montre quelles cases sont revisitées quand la bande s’élargit.}
